% ==============================================================================
% PREFACIO: DECLARACIÓN SOBRE EL USO DE INTELIGENCIA ARTIFICIAL
% ==============================================================================
% Adaptado (2026-10-12) de la estructura de la declaración del otro integrante:
% rol del autor, herramientas empleadas, y cierre de responsabilidad.
%
% Este tomo todavía no tiene agradecimientos. Cuando existan, van antes de la
% declaración (esta es el último bloque del prefacio).
%
% En main.tex, el \chapter*{Prefacio} lo pone el documento principal; en modo
% independiente se agrega aquí.
\ifdefined\modoensamble
\else
	\input{0-preambulo_standalone}
	\chapter*{Prefacio}
\fi

\subsection*{Declaración sobre el uso de inteligencia artificial}

El autor definió el enfoque técnico de su componente, seleccionó los vectores de ataque y su
mapeo a OWASP Top 10 y MITRE ATT\&CK, auditó y etiquetó los datos, construyó el entorno de
evaluación, diseñó la capa de captura y respuesta (respuesta HTTP 403, registro de eventos,
\textit{webhooks} e interfaz de línea de comandos), calibró el umbral de bloqueo, elaboró el
corpus de evaluación, ejecutó las pruebas de penetración, las mediciones de latencia, memoria y CPU
y la validación en DVNA, node-api-goat y dvws-node, e implementó el extractor inicial de 72
características y el paso a \textit{worker\_threads}. Decidió además el alcance y los
resultados que se reportan.

Para la búsqueda bibliográfica empleó Gemini (Google) como herramienta de apoyo para identificar
fuentes y temas potencialmente relevantes; leyó las fuentes originales y conservó únicamente las que
respaldaban lo que se afirma en el texto. Para la documentación técnica, la redacción y la
verificación de coherencia de cifras, criterios y atribución de roles entre este trabajo y el del
componente de Inteligencia Artificial empleó a Claude, de Anthropic, mediante conversación, y a
Claude Code, mediante un agente de código, que además se usó para refactorizar código y como
revisor adicional de \textit{pull requests}; esa revisión complementó, y no sustituyó, la revisión
cruzada obligatoria entre los integrantes del proyecto descrita en la
Sección~\ref{subsec:met-flujo}. Ninguna de estas herramientas se utilizó para generar datos,
mediciones ni resultados: las cifras provienen de los archivos y las ejecuciones de los
repositorios del proyecto, y las derivadas de ellos (medianas, intervalos de confianza) se
calcularon a partir de esos archivos. El uso de estas herramientas se declara conforme a la \textit{Guía para el uso ético y responsable de la inteligencia artificial} de la Universidad del Valle de Guatemala. El autor dirigió ese trabajo, contrastó las cifras con sus
fuentes, revisó las referencias bibliográficas contra las fuentes originales, salvo las señaladas
en la bibliografía como no accesibles, puede explicar el contenido y asume la responsabilidad del
documento.

\ifdefined\modoensamble
\else
	\end{document}
\fi
